% !TEX program = xelatex
\documentclass[aspectratio=169,11pt]{beamer}
\usetheme{Madrid}
\usepackage{fontspec}
\setsansfont{DejaVu Sans}
\setmainfont{DejaVu Serif}
\usepackage{polyglossia}
\setdefaultlanguage{turkish}
\usepackage{amsmath,amssymb}
\usepackage{booktabs}
\usepackage{multicol}
\usepackage{xcolor}

% ---------- Renkler ----------
\definecolor{anamavi}{RGB}{0,45,98}
\definecolor{cozumyesil}{RGB}{0,110,50}
\definecolor{uyarikirmizi}{RGB}{180,20,20}
\setbeamercolor{structure}{fg=anamavi}
\setbeamercolor{frametitle}{bg=anamavi,fg=white}
\setbeamercolor{title}{bg=anamavi,fg=white}
\setbeamercolor{palette primary}{bg=anamavi,fg=white}
\setbeamercolor{palette secondary}{bg=anamavi!85!black,fg=white}
\setbeamercolor{palette tertiary}{bg=anamavi!70!black,fg=white}
\setbeamercolor{block title}{bg=anamavi,fg=white}
\setbeamercolor{block body}{bg=anamavi!7}
\setbeamercolor{block title example}{bg=cozumyesil,fg=white}
\setbeamercolor{block body example}{bg=cozumyesil!8}
\setbeamercolor{block title alerted}{bg=uyarikirmizi,fg=white}
\setbeamercolor{block body alerted}{bg=uyarikirmizi!7}
\setbeamertemplate{navigation symbols}{}
\setbeamertemplate{blocks}[rounded][shadow=false]

% ---------- Soru / Çözüm kalıbı ----------
% \soru{başlık}{soru metni}{çözüm metni}: önce soru, bir tıklamayla çözüm
\newcommand{\soru}[3]{%
\begin{frame}{#1}
  \begin{block}{Soru}\large #2\end{block}
  \pause
  \begin{exampleblock}{Çözüm}\large #3\end{exampleblock}
\end{frame}}
\newcommand{\ab}[1]{\overline{#1}}

\title[MMAT105 -- Hafta 5]{MMAT105 Genel Matematik}
\subtitle{Hafta 5: Birinci Dereceden Bir Bilinmeyenli Denklemler}
\author{}
\date{2026--2027 Güz Dönemi}

\begin{document}

% =====================================================
\begin{frame}
\titlepage
\end{frame}

% (iletişim çerçevesi kaldırıldı)

\begin{frame}{Geçen Hafta Hatırlatma}
\begin{columns}[t]
\column{0.56\textwidth}
\textbf{Hafta 4: Rasyonel ve ondalık sayılar}
\begin{enumerate}
\item \(\dfrac{2}{3}+\dfrac{5}{6}-\dfrac{1}{4}=?\)\\[4pt]
\item \(0{,}\overline{3}+0{,}1\overline{6}=?\)\\[4pt]
\item \(a=\dfrac{2}{3},\ b=\dfrac{3}{5},\ c=\dfrac{5}{8}\) sayılarını küçükten büyüğe sıralayınız.
\end{enumerate}
\column{0.40\textwidth}
\pause
\textbf{Cevaplar:}
\begin{enumerate}
\item \(\dfrac{8+10-3}{12}=\dfrac{15}{12}=\dfrac{5}{4}\)\\[4pt]
\item \(\dfrac{1}{3}+\dfrac{15}{90}=\dfrac{1}{3}+\dfrac{1}{6}=\dfrac{1}{2}\)\\[4pt]
\item \(b<c<a\) \quad (\(0{,}6<0{,}625<0{,}\overline{6}\))
\end{enumerate}
\end{columns}
\end{frame}

\begin{frame}{Bugün}
\begin{columns}[t]
\column{0.32\textwidth}
\begin{block}{Ders 1}
\begin{itemize}
\item Denklem, eşitlik, özdeşlik
\item Eşitliğin özellikleri
\item \(ax+b=0\) ve çözümü
\item Parantezli denklemler
\item Sağlama
\end{itemize}
\end{block}
\column{0.32\textwidth}
\begin{block}{Ders 2}
\begin{itemize}
\item Kesirli denklemler
\item Ondalıklı katsayılar
\item Özel durumlar: \(\mathbb{R}\), \(\emptyset\)
\item Parametreli sorular
\item Paydada bilinmeyen
\end{itemize}
\end{block}
\column{0.32\textwidth}
\begin{block}{Ders 3}
\begin{itemize}
\item Sözel ifadeyi denkleme çevirme
\item Sayı problemleri
\item Ardışık sayılar
\item Para, yaş, yol: giriş
\end{itemize}
\end{block}
\end{columns}
\end{frame}

% =====================================================
\section{Ders 1: Denklem Kavramı ve Çözümü}
\begin{frame}
\centering\Large\color{anamavi}\textbf{Ders 1}\\[6pt]\normalsize Denklem, Eşitlik, Özdeşlik; \(ax+b=0\) Denkleminin Çözümü
\end{frame}

\subsection{Denklem, eşitlik, özdeşlik}
\begin{frame}{Denklem, Eşitlik, Özdeşlik}
\begin{columns}[c]
\column{0.64\textwidth}
\begin{block}{Tanımlar}
\begin{itemize}
\item \textbf{Eşitlik:} İki ifadenin eşit olduğunu bildiren yazı: \(3+5=8\).
\item \textbf{Denklem:} Bilinmeyen içeren, her değer için doğru \textbf{olmayan} eşitlik: \(2x+3=11\) (yalnız \(x=4\) için doğru); hiçbir değer için doğru olmayabilir de.
\item \textbf{Özdeşlik:} Bilinmeyenin \textbf{her} değeri için doğru olan eşitlik: \(2(x+3)=2x+6\).
\item Denklemi sağlayan değere \textbf{kök}, köklerin kümesine \textbf{çözüm kümesi} (\(\text{ÇK}\)) denir.
\end{itemize}
\end{block}
\column{0.33\textwidth}
\centering
\setlength{\unitlength}{1mm}
\begin{picture}(40,30)
\thicklines
\put(20,0){\line(0,1){20}}
\put(12,0){\line(1,0){16}}
\put(2,20){\line(1,0){36}}
\put(2,20){\line(0,-1){6}}\put(38,20){\line(0,-1){6}}
\put(-3,10){\framebox(10,4){\scriptsize \(3x+2\)}}
\put(33,10){\framebox(10,4){\scriptsize \(x+10\)}}
\put(20,20){\circle*{1.5}}
\end{picture}\\[2pt]
{\small Terazi dengede \(\Leftrightarrow\) iki kefe eşit}
\end{columns}
\end{frame}

\begin{frame}{Eşitliğin Özellikleri ve \(ax+b=0\)}
\begin{block}{Terazi kuralı: iki kefeye aynı işlemi yap}
\(a=b\) ise \quad \(a+c=b+c\), \quad \(a-c=b-c\), \quad \(a\cdot c=b\cdot c\), \quad \(c\neq 0\) ise \(\dfrac{a}{c}=\dfrac{b}{c}\).
\end{block}
\begin{block}{Birinci dereceden bir bilinmeyenli denklem}
\[ax+b=0\quad (a\neq 0)\qquad\Longrightarrow\qquad x=-\frac{b}{a},\qquad \text{ÇK}=\left\{-\frac{b}{a}\right\}\]
Adımlar: parantezleri aç \(\to\) bilinmeyenleri bir tarafa, sayıları öbür tarafa topla \(\to\) katsayıya böl \(\to\) \textbf{sağla}.
\end{block}
\begin{alertblock}{Sık yapılan hatalar}
Karşı tarafa geçen terimin işareti değişir: \(x+5=9\Rightarrow x=9-5\). \quad \(-(x-3)=-x+3\). \quad \textbf{Sıfıra bölme yoktur.}
\end{alertblock}
\end{frame}

\soru{Kavramlar -- Örnek 1}
{Aşağıdakilerden hangisi denklem, hangisi özdeşlik, hangisi sayısal eşitliktir?
\[\text{(a) } 2x+3=11\qquad \text{(b) } 2(x+3)=2x+6\qquad \text{(c) } 3+5=8\]}
{(a) \textbf{Denklem}: yalnız \(x=4\) için doğru (\(8+3=11\)).\\[3pt]
(b) \textbf{Özdeşlik}: parantez açılınca \(2x+6=2x+6\); her \(x\) için doğru.\\[3pt]
(c) \textbf{Sayısal eşitlik}: bilinmeyen yok.}

\soru{Terazi -- Örnek 2}
{Dengedeki bir terazinin bir kefesinde \(3\) özdeş kutu ve \(2\) kg, diğer kefesinde \(1\) kutu ve \(10\) kg vardır. Bir kutu kaç kg'dır?}
{Kutu \(x\) kg olsun: \quad \(3x+2=x+10\)\\[3pt]
İki kefeden \(1\) kutu ve \(2\) kg çıkaralım: \quad \(2x=8\ \Rightarrow\ x=4\) kg.\\[3pt]
Sağlama: \(3\cdot 4+2=14=4+10\).}

\soru{Temel Çözüm -- Örnek 3}
{\(5x-7=18\) denkleminin çözüm kümesini bulunuz.}
{\[5x=18+7=25\ \Rightarrow\ x=5\]
Sağlama: \(5\cdot 5-7=18\) \checkmark \qquad \(\text{ÇK}=\{5\}\)}

\soru{Temel Çözüm -- Örnek 4}
{\(3x+5=x+9\) denklemini çözünüz.}
{Bilinmeyenler solda, sayılar sağda toplanır:
\[3x-x=9-5\ \Rightarrow\ 2x=4\ \Rightarrow\ x=2\]
Sağlama: \(3\cdot 2+5=11=2+9\) \checkmark \qquad \(\text{ÇK}=\{2\}\)}

\soru{Temel Çözüm -- Örnek 5}
{\(7-3x=2x-8\) denklemini çözünüz.}
{\[7+8=2x+3x\ \Rightarrow\ 15=5x\ \Rightarrow\ x=3\]
Sağlama: \(7-9=-2\), \ \(6-8=-2\) \checkmark \qquad \(\text{ÇK}=\{3\}\)\\[3pt]
İpucu: \(x\)'leri katsayısı büyük olan tarafa toplamak işaret hatasını azaltır.}

\subsection{Parantezli denklemler}
\soru{Parantezli -- Örnek 6}
{\(3(x-2)-2(x+1)=4\) denklemini çözünüz.}
{Parantezleri açalım (dikkat: \(-2\cdot(+1)=-2\)):
\[3x-6-2x-2=4\ \Rightarrow\ x-8=4\ \Rightarrow\ x=12\]
Sağlama: \(3\cdot 10-2\cdot 13=30-26=4\) \checkmark}

\soru{Parantezli -- Örnek 7}
{\(5-2(3-x)=4(x-1)+1\) denklemini çözünüz.}
{\[5-6+2x=4x-4+1\ \Rightarrow\ 2x-1=4x-3\ \Rightarrow\ 2=2x\ \Rightarrow\ x=1\]
Sağlama: \(5-2\cdot 2=1\), \ \(4\cdot 0+1=1\) \checkmark \qquad \(\text{ÇK}=\{1\}\)}

\soru{İç İçe Parantez -- Örnek 8}
{\(2\big[x-3(x-1)\big]=x+1\) denklemini çözünüz.}
{En içten başlayalım: \(x-3(x-1)=x-3x+3=-2x+3\)
\[2(-2x+3)=x+1\ \Rightarrow\ -4x+6=x+1\ \Rightarrow\ 5=5x\ \Rightarrow\ x=1\]
Sağlama: \(2[1-0]=2=1+1\) \checkmark}

\soru{Sağlama -- Örnek 9}
{Bir öğrenci \(4(x+1)=2x-6\) denkleminin kökünü \(x=-1\) bulmuştur. Sağlama yaparak kontrol ediniz; yanlışsa doğru kökü bulunuz.}
{\(x=-1\): \ sol \(=4\cdot 0=0\), \ sağ \(=-2-6=-8\). \ \(0\neq-8\) \(\Rightarrow\) \textbf{yanlış}.\\[3pt]
Doğrusu: \(4x+4=2x-6\ \Rightarrow\ 2x=-10\ \Rightarrow\ x=-5\)\\[3pt]
Sağlama: \(4\cdot(-4)=-16\), \ \(2\cdot(-5)-6=-16\) \checkmark}

\soru{Sıfıra Bölme Tuzağı -- Örnek 10}
{Bir öğrenci \(3x=5x\) denkleminde iki tarafı \(x\)'e bölüp \(3=5\) bulmuş ve ``çözüm yok'' demiştir. Hata nerededir?}
{\(x=0\) olabileceği için \(x\)'e bölmek \textbf{sıfıra bölme} riski taşır.\\[3pt]
Doğrusu: \(5x-3x=0\ \Rightarrow\ 2x=0\ \Rightarrow\ x=0\) \qquad \(\text{ÇK}=\{0\}\)\\[3pt]
Kural: Bilinmeyen içeren bir ifadeye bölmeyin; terimleri bir tarafa toplayın.}

\soru{Karışık -- Örnek 11}
{\(-(x-4)+3=2(5-x)\) denklemini çözünüz.}
{\[-x+4+3=10-2x\ \Rightarrow\ -x+7=10-2x\ \Rightarrow\ x=3\]
Sağlama: \(-(-1)+3=4\), \ \(2\cdot 2=4\) \checkmark \qquad \(\text{ÇK}=\{3\}\)}

\begin{frame}{Ders 1 -- Sınıfta Siz Çözün}
\begin{columns}[t]
\column{0.58\textwidth}
\begin{enumerate}
\item \(4x-9=11\)
\item \(2(x+3)=x+10\)
\item \(7-(x-2)=3x-3\) \(\star\)
\item \(3(2x-1)-2(x+4)=5\) \(\star\)
\item \(5x=-2x\) \(\star\)
\end{enumerate}
\column{0.36\textwidth}
\pause
\textbf{Cevaplar:}
\begin{enumerate}
\item \(x=5\)
\item \(x=4\)
\item \(x=3\)
\item \(x=4\)
\item \(x=0\) \quad (\(7x=0\))
\end{enumerate}
\end{columns}
\vspace{4pt}{\footnotesize \(\star\) işaretli soruların adım adım çözümü sonraki slaytlardadır.}
\end{frame}

% ---- Ders 1 alıştırma çözümleri ----
\soru{Ders 1 -- Alıştırma Çözümü (Soru 3)}
{\(7-(x-2)=3x-3\) denklemini çözünüz.}
{Parantezin önündeki eksi, içteki her terimin işaretini değiştirir:
\[7-x+2=3x-3\ \Rightarrow\ 9-x=3x-3\]
\[9+3=3x+x\ \Rightarrow\ 12=4x\ \Rightarrow\ x=3\]
Sağlama: \(7-1=6\), \ \(9-3=6\) \checkmark \qquad \(\text{ÇK}=\{3\}\)}

\soru{Ders 1 -- Alıştırma Çözümü (Soru 4)}
{\(3(2x-1)-2(x+4)=5\) denklemini çözünüz.}
{Dağılma: \(6x-3-2x-8=5\) \quad (\(-2\cdot 4=-8\))
\[4x-11=5\ \Rightarrow\ 4x=16\ \Rightarrow\ x=4\]
Sağlama: \(3\cdot 7-2\cdot 8=21-16=5\) \checkmark \qquad \(\text{ÇK}=\{4\}\)}

\soru{Ders 1 -- Alıştırma Çözümü (Soru 5)}
{\(5x=-2x\) denklemini çözünüz.}
{Bilinmeyenler bir tarafa toplanır: \(5x+2x=0\ \Rightarrow\ 7x=0\ \Rightarrow\ x=0\).\\[3pt]
\textbf{Dikkat:} İki tarafı \(x\)'e bölmek yanlıştır (\(x=0\) olabilir; \(0\)'a bölme tanımsız). Bölünürse \(5=-2\) gibi yanlış bir sonuç çıkar.\\[3pt]
Sağlama: \(5\cdot 0=0=-2\cdot 0\) \checkmark \qquad \(\text{ÇK}=\{0\}\)}

% =====================================================
\section{Ders 2: Kesirli Denklemler ve Özel Durumlar}
\begin{frame}
\centering\Large\color{anamavi}\textbf{Ders 2}\\[6pt]\normalsize Kesirli ve Ondalıklı Denklemler, Özel Durumlar, Parametreler, Paydada Bilinmeyen
\end{frame}

\subsection{Kesirli ve ondalıklı denklemler}
\begin{frame}{Kesirli ve Ondalıklı Katsayılı Denklemler}
\begin{block}{Kural}
\begin{itemize}
\item \textbf{Kesirli:} Her iki tarafın \textbf{her terimini} paydaların EKOK'u ile çarp; paydalar gider.
\item \textbf{Ondalıklı:} En çok \(k\) ondalık basamak varsa her terimi \(10^{k}\) ile çarp.
\end{itemize}
\end{block}
\begin{alertblock}{Sık yapılan hatalar}
\begin{itemize}
\item Paydası olmayan terimi çarpmayı unutmak: \(\dfrac{x}{2}+1=3\ \xrightarrow{\cdot 2}\ x+2=6\) \ (\(x+1=6\) değil).
\item Kesir çizgisi parantezdir: \(-\dfrac{x-3}{2}\ \xrightarrow{\cdot 2}\ -(x-3)=-x+3\).
\end{itemize}
\end{alertblock}
\end{frame}

\soru{Kesirli -- Örnek 12}
{\(\dfrac{x}{2}+\dfrac{x}{3}=10\) denklemini çözünüz.}
{\(\text{EKOK}(2,3)=6\); her terimi \(6\) ile çarpalım:
\[3x+2x=60\ \Rightarrow\ 5x=60\ \Rightarrow\ x=12\]
Sağlama: \(6+4=10\) \checkmark}

\soru{Kesirli -- Örnek 13}
{\(\dfrac{x-1}{3}-\dfrac{x+2}{4}=1\) denklemini çözünüz.}
{\(\text{EKOK}(3,4)=12\):
\[4(x-1)-3(x+2)=12\ \Rightarrow\ 4x-4-3x-6=12\ \Rightarrow\ x=22\]
Sağlama: \(\dfrac{21}{3}-\dfrac{24}{4}=7-6=1\) \checkmark}

\soru{Kesirli -- Örnek 14}
{\(\dfrac{2x+1}{5}-\dfrac{x-3}{2}=1\) denklemini çözünüz.}
{\(\text{EKOK}(5,2)=10\):
\[2(2x+1)-5(x-3)=10\ \Rightarrow\ 4x+2-5x+15=10\ \Rightarrow\ -x=-7\ \Rightarrow\ x=7\]
Sağlama: \(\dfrac{15}{5}-\dfrac{4}{2}=3-2=1\) \checkmark}

\soru{Ondalıklı -- Örnek 15}
{\(0{,}2x+1{,}5=0{,}5x-0{,}3\) denklemini çözünüz.}
{En çok \(1\) ondalık basamak var \(\Rightarrow\) \(10\) ile çarpalım:
\[2x+15=5x-3\ \Rightarrow\ 18=3x\ \Rightarrow\ x=6\]
Sağlama: \(1{,}2+1{,}5=2{,}7\), \ \(3-0{,}3=2{,}7\) \checkmark}

\soru{Ondalıklı -- Örnek 16}
{\(0{,}12x+0{,}3=0{,}06x+0{,}9\) denklemini çözünüz.}
{En çok \(2\) ondalık basamak var \(\Rightarrow\) \(100\) ile çarpalım:
\[12x+30=6x+90\ \Rightarrow\ 6x=60\ \Rightarrow\ x=10\]
Sağlama: \(1{,}2+0{,}3=1{,}5\), \ \(0{,}6+0{,}9=1{,}5\) \checkmark}

\subsection{Özel durumlar ve parametreler}
\begin{frame}{Özel Durumlar: Sonsuz Çözüm ve Çözümsüzlük}
\begin{block}{\(ax+b=0\) denkleminde}
\begin{tabular}{@{}lll@{}}
\(a\neq 0\) & tek kök \(x=-\dfrac{b}{a}\) & \(\text{ÇK}=\left\{-\dfrac{b}{a}\right\}\)\\[8pt]
\(a=0,\ b=0\) & \(0\cdot x=0\): her \(x\) sağlar (özdeşlik) & \(\text{ÇK}=\mathbb{R}\)\\[4pt]
\(a=0,\ b\neq 0\) & \(0\cdot x=-b\neq 0\): hiçbir \(x\) sağlamaz & \(\text{ÇK}=\emptyset\)
\end{tabular}
\end{block}
\begin{block}{\(ax+b=cx+d\) biçiminde}
\(\text{ÇK}=\mathbb{R}\ \Leftrightarrow\ a=c\) ve \(b=d\); \qquad \(\text{ÇK}=\emptyset\ \Leftrightarrow\ a=c\) ve \(b\neq d\).
\end{block}
\begin{alertblock}{Dikkat}
\(\emptyset\) ile \(\{0\}\) farklıdır: \(\{0\}\) kümesinin bir elemanı (\(0\)) vardır.
\end{alertblock}
\end{frame}

\soru{Özel Durum -- Örnek 17}
{Çözüm kümelerini bulunuz:\\[3pt] (a) \(2(x+3)-x=x+6\) \qquad\qquad (b) \(3x-2=3(x+1)\)}
{(a) \(2x+6-x=x+6\ \Rightarrow\ x+6=x+6\ \Rightarrow\ 0\cdot x=0\) \quad \(\text{ÇK}=\mathbb{R}\)\\[4pt]
(b) \(3x-2=3x+3\ \Rightarrow\ 0\cdot x=5\) \quad (hiçbir \(x\) sağlamaz) \quad \(\text{ÇK}=\emptyset\)}

\soru{Parametre -- Örnek 18}
{\((m-2)x+4=3x-n\) denkleminin çözüm kümesi \(\mathbb{R}\) ise \(m+n\) kaçtır?}
{Her \(x\) için doğru olmalı: \(x\)'in katsayıları ve sabitler eşit olmalı.
\[m-2=3\ \Rightarrow\ m=5,\qquad 4=-n\ \Rightarrow\ n=-4\]
\[m+n=5+(-4)=1\]}

\soru{Parametre -- Örnek 19}
{\((a+1)x-3=2x+b\) denkleminin çözüm kümesi boş küme ise \(a\) kaçtır? \(b\) için hangi koşul gerekir?}
{Çözüm kümesi boş \(\Rightarrow\) katsayılar eşit, sabitler farklı:
\[a+1=2\ \Rightarrow\ a=1,\qquad -3\neq b\ \Rightarrow\ b\neq-3\]
(\(b=-3\) olsaydı \(\text{ÇK}=\mathbb{R}\) olurdu.)}

\soru{Kök Verilince -- Örnek 20}
{\(ax-4=2x+a\) denkleminin kökü \(3\) ise \(a\) kaçtır?}
{Kök denklemi sağlar: \(x=3\) yazalım.
\[3a-4=6+a\ \Rightarrow\ 2a=10\ \Rightarrow\ a=5\]
Kontrol: \(5x-4=2x+5\Rightarrow 3x=9\Rightarrow x=3\) \checkmark}

\subsection{Paydada bilinmeyen}
\begin{frame}{Paydada Bilinmeyen Olan Denklemler}
\begin{block}{Kural}
\begin{enumerate}
\item Önce \textbf{tanımsızlık koşulunu} yaz: payda \(\neq 0\) \ (ör. \(x-2\neq 0\Rightarrow x\neq 2\)).
\item Paydaları yok et (paydaların EKOK'u ile çarp), denklemi çöz.
\item Bulunan değer koşula uymuyorsa kök \textbf{değildir}; atılır.
\end{enumerate}
\end{block}
\begin{alertblock}{Sık yapılan hatalar}
\(\dfrac{x+1}{x-3}=\dfrac{4}{x-3}\) denkleminden \(x=3\) bulup ``\(\text{ÇK}=\{3\}\)'' yazmak. \ \(x=3\) paydayı sıfır yapar!
\end{alertblock}
\end{frame}

\soru{Paydada Bilinmeyen -- Örnek 21}
{\(\dfrac{6}{x-1}=3\) denklemini çözünüz.}
{Koşul: \(x\neq 1\).
\[6=3(x-1)\ \Rightarrow\ 6=3x-3\ \Rightarrow\ x=3\]
\(3\neq 1\), koşula uyar. Sağlama: \(\dfrac{6}{2}=3\) \checkmark \qquad \(\text{ÇK}=\{3\}\)}

\soru{Paydada Bilinmeyen -- Örnek 22}
{\(\dfrac{x+3}{x-2}=2\) denklemini çözünüz.}
{Koşul: \(x\neq 2\).
\[x+3=2(x-2)\ \Rightarrow\ x+3=2x-4\ \Rightarrow\ x=7\]
Sağlama: \(\dfrac{10}{5}=2\) \checkmark \qquad \(\text{ÇK}=\{7\}\)}

\soru{Tuzak -- Örnek 23}
{\(\dfrac{x+1}{x-3}=\dfrac{4}{x-3}\) denkleminin çözüm kümesini bulunuz.}
{Koşul: \(x\neq 3\). \ Paydalar aynı: \(x+1=4\ \Rightarrow\ x=3\).\\[3pt]
Ama \(x=3\) koşulu bozar (payda \(0\) olur) \(\Rightarrow\) kök değildir.
\[\text{ÇK}=\emptyset\]}

\soru{Paydada Bilinmeyen -- Örnek 24}
{\(\dfrac{2}{x}+\dfrac{1}{2x}=\dfrac{5}{4}\) denklemini çözünüz.}
{Koşul: \(x\neq 0\). \ \(\text{EKOK}(x,\,2x,\,4)=4x\):
\[8+2=5x\ \Rightarrow\ x=2\]
Sağlama: \(1+\dfrac{1}{4}=\dfrac{5}{4}\) \checkmark \qquad \(\text{ÇK}=\{2\}\)}

\begin{frame}{Ders 2 -- Sınıfta Siz Çözün}
\begin{columns}[t]
\column{0.62\textwidth}
\begin{enumerate}
\item \(\dfrac{x}{3}+\dfrac{x}{4}=7\)
\item \(\dfrac{x+2}{3}-\dfrac{x-1}{2}=1\) \(\star\)
\item \(0{,}3x-0{,}4=0{,}1x+0{,}6\)
\item \((a-3)x+2=5x+b\) denkleminin çözüm kümesi \(\mathbb{R}\) ise \(a+b=?\) \(\star\)
\item \(\dfrac{2x+5}{x+1}=3\) \(\star\)
\end{enumerate}
\column{0.34\textwidth}
\pause
\textbf{Cevaplar:}
\begin{enumerate}
\item \(x=12\)\\[6pt]
\item \(x=1\)\\[6pt]
\item \(x=5\)
\item \(a=8,\ b=2\Rightarrow 10\)\\[6pt]
\item \(x=2\) \ (\(x\neq-1\))
\end{enumerate}
\end{columns}
\vspace{4pt}{\footnotesize \(\star\) işaretli soruların adım adım çözümü sonraki slaytlardadır.}
\end{frame}

% ---- Ders 2 alıştırma çözümleri ----
\soru{Ders 2 -- Alıştırma Çözümü (Soru 2)}
{\(\dfrac{x+2}{3}-\dfrac{x-1}{2}=1\) denklemini çözünüz.}
{Her terim \(\text{EKOK}(3,2)=6\) ile çarpılır (sağdaki \(1\) dahil!):
\[2(x+2)-3(x-1)=6\ \Rightarrow\ 2x+4-3x+3=6\]
\[-x+7=6\ \Rightarrow\ x=1\]
Sağlama: \(\dfrac{3}{3}-\dfrac{0}{2}=1\) \checkmark \qquad \(\text{ÇK}=\{1\}\)}

\soru{Ders 2 -- Alıştırma Çözümü (Soru 4)}
{\((a-3)x+2=5x+b\) denkleminin çözüm kümesi \(\mathbb{R}\) ise \(a+b\) kaçtır?}
{Çözüm kümesi \(\mathbb{R}\) ise iki taraf özdeştir: \(x\)'in katsayıları ve sabitler eşit olmalı.\\[3pt]
\(x\)'in katsayıları: \(a-3=5\ \Rightarrow\ a=8\)\\[3pt]
Sabitler: \(2=b\ \Rightarrow\ b=2\)
\[a+b=8+2=10\]}

\soru{Ders 2 -- Alıştırma Çözümü (Soru 5)}
{\(\dfrac{2x+5}{x+1}=3\) denklemini çözünüz.}
{Koşul: \(x+1\neq 0\ \Rightarrow\ x\neq -1\).\\[3pt]
İki taraf \((x+1)\) ile çarpılır:
\[2x+5=3(x+1)=3x+3\ \Rightarrow\ x=2\]
\(2\neq-1\) \checkmark \quad Sağlama: \(\dfrac{9}{3}=3\) \checkmark \qquad \(\text{ÇK}=\{2\}\)}

% =====================================================
\section{Ders 3: Sözel İfadeler ve Problemler}
\begin{frame}
\centering\Large\color{anamavi}\textbf{Ders 3}\\[6pt]\normalsize Sözel İfadeleri Denkleme Çevirme, Sayı ve Ardışık Sayı Problemleri, Para -- Yaş -- Yol (Giriş)
\end{frame}

\subsection{Sözel ifadeden denkleme}
\begin{frame}{Sözel İfadeden Denkleme}
\begin{columns}[t]
\column{0.52\textwidth}
\begin{block}{Çeviri sözlüğü (sayı: \(x\))}
\small
\begin{tabular}{@{}ll@{}}
3 katı & \(3x\)\\
5 fazlası / 5 eksiği & \(x+5\) / \(x-5\)\\
yarısı, üçte biri & \(\dfrac{x}{2}\), \ \(\dfrac{x}{3}\)\\[4pt]
3 katının 5 fazlası & \(3x+5\)\\
5 fazlasının 3 katı & \(3(x+5)\)\\
ardışık tam sayılar & \(x,\ x+1,\ x+2\)\\
ardışık çift / tek & \(x,\ x+2,\ x+4\)
\end{tabular}
\end{block}
\column{0.44\textwidth}
\begin{block}{Problem çözme adımları}
\begin{enumerate}
\item Bilinmeyeni seç (\(x\) = ?)
\item Diğer nicelikleri \(x\) ile yaz
\item Denklemi kur ve çöz
\item Sağla, \textbf{sorulanı} cevapla
\end{enumerate}
\end{block}
\begin{alertblock}{Sık yapılan hatalar}
\small ``sayının \(5\) eksiği'' \(x-5\), ama ``\(5\)'in sayı kadar eksiği'' \(5-x\).\\
Son adımda sorulan \(x\) değil, \(x+2\) olabilir!
\end{alertblock}
\end{columns}
\end{frame}

\soru{Çeviri -- Örnek 25}
{Cebirsel ifadeye çeviriniz (sayı \(x\)):\\
(a) bir sayının \(3\) katının \(5\) fazlası \quad (b) bir sayının \(5\) fazlasının \(3\) katı\\
(c) bir sayının yarısının \(4\) eksiği \quad (d) bir sayının \(4\) eksiğinin yarısı}
{\[\text{(a) } 3x+5\qquad \text{(b) } 3(x+5)\qquad \text{(c) } \frac{x}{2}-4\qquad \text{(d) } \frac{x-4}{2}\]
Son söylenen işlem en son yapılır: ``\dots\ fazlasının \(3\) katı'' \(\Rightarrow\) önce ekle, sonra çarp.}

\subsection{Sayı problemleri}
\soru{Sayı Problemi -- Örnek 26}
{Bir sayının \(3\) katının \(5\) fazlası \(29\) ise bu sayı kaçtır?}
{\[3x+5=29\ \Rightarrow\ 3x=24\ \Rightarrow\ x=8\]
Sağlama: \(3\cdot 8+5=29\) \checkmark}

\soru{Sayı Problemi -- Örnek 27}
{Bir sayının \(4\) fazlasının \(3\) katı, aynı sayının \(5\) katından \(2\) eksiktir. Bu sayı kaçtır?}
{\[3(x+4)=5x-2\ \Rightarrow\ 3x+12=5x-2\ \Rightarrow\ 14=2x\ \Rightarrow\ x=7\]
Sağlama: \(3\cdot 11=33\), \ \(35-2=33\) \checkmark}

\soru{Sayı Problemi -- Örnek 28}
{Bir sayının yarısı ile üçte birinin toplamı \(25\) ise bu sayı kaçtır?}
{\[\frac{x}{2}+\frac{x}{3}=25\ \xrightarrow{\cdot 6}\ 3x+2x=150\ \Rightarrow\ x=30\]
Sağlama: \(15+10=25\) \checkmark}

\soru{Bilmece -- Örnek 29}
{Havada uçan kazlar yerdeki topal kaza: ``Biz, bizim kadar, bizim yarımız kadar, yarımızın yarısı kadar ve bir de sen olursan ancak \(100\) oluruz.'' Havada kaç kaz vardır?}
{\[x+x+\frac{x}{2}+\frac{x}{4}+1=100\ \xrightarrow{\cdot 4}\ 4x+4x+2x+x+4=400\]
\[11x=396\ \Rightarrow\ x=36\qquad(36+36+18+9+1=100\ \checkmark)\]}

\subsection{Ardışık sayılar}
\soru{Ardışık -- Örnek 30}
{Ardışık üç tam sayının toplamı \(96\) ise en büyüğü kaçtır?}
{Sayılar \(x,\ x+1,\ x+2\):
\[3x+3=96\ \Rightarrow\ x=31\]
Sayılar \(31,\ 32,\ 33\) \(\Rightarrow\) en büyük: \(33\).}

\soru{Ardışık -- Örnek 31}
{Ardışık üç çift sayıdan en küçüğünün \(3\) katı, en büyüğünün \(2\) katından \(14\) fazladır. Bu üç sayının toplamı kaçtır?}
{Sayılar \(x,\ x+2,\ x+4\):
\[3x=2(x+4)+14\ \Rightarrow\ 3x=2x+22\ \Rightarrow\ x=22\]
Sayılar \(22,\ 24,\ 26\); toplam \(72\). \ (Sağlama: \(66=52+14\) \checkmark)}

\subsection{Para, yaş, yol}
\soru{Para -- Örnek 32}
{Ali'nin parası Ayşe'nin parasının \(3\) katıdır. Ali, Ayşe'ye \(40\) TL verirse paraları eşit oluyor. Ayşe'nin parası kaç TL'dir?}
{Ayşe: \(x\), \ Ali: \(3x\).
\[3x-40=x+40\ \Rightarrow\ 2x=80\ \Rightarrow\ x=40\]
Ayşe \(40\) TL, Ali \(120\) TL. \ (Sağlama: \(80=80\) \checkmark)}

\soru{Yaş -- Örnek 33}
{Bir annenin yaşı kızının yaşının \(4\) katıdır. \(6\) yıl sonra annenin yaşı kızının yaşının \(3\) katı olacaktır. Kızın bugünkü yaşı kaçtır?}
{Kız: \(x\), anne: \(4x\). \ \(6\) yıl sonra herkes \(6\) yaş büyür:
\[4x+6=3(x+6)\ \Rightarrow\ 4x+6=3x+18\ \Rightarrow\ x=12\]
Kız \(12\), anne \(48\). \ (Sağlama: \(54=3\cdot 18\) \checkmark)}

\soru{Yol -- Örnek 34}
{Aralarında \(450\) km olan iki şehirden iki araç aynı anda birbirine doğru yola çıkıyor. Hızları \(60\) km/sa ve \(90\) km/sa ise kaç saat sonra karşılaşırlar?}
{Yol \(=\) hız \(\times\) zaman. \(t\) saatte alınan yolların toplamı \(450\) km:
\[60t+90t=450\ \Rightarrow\ 150t=450\ \Rightarrow\ t=3 \text{ saat}\]}

\soru{Şifre -- Örnek 35}
{Dört basamaklı bir PIN kodunda soldan ilk iki basamağın oluşturduğu sayı, son iki basamağın oluşturduğu sayının \(2\) katından \(1\) fazladır; bu iki sayının toplamı \(46\)'dır. PIN kodunun rakamları toplamı kaçtır?}
{Son iki basamak \(x\), ilk iki basamak \(2x+1\):
\[(2x+1)+x=46\ \Rightarrow\ 3x=45\ \Rightarrow\ x=15\]
İlk iki basamak \(31\) \(\Rightarrow\) PIN: \(3115\); rakamlar toplamı \(3+1+1+5=10\).}

\begin{frame}{Ders 3 -- Sınıfta Siz Çözün}
\begin{columns}[t]
\column{0.64\textwidth}
\begin{enumerate}
\item Bir sayının \(2\) katının \(7\) eksiği \(25\) ise sayı kaçtır?
\item Ardışık dört tek sayının toplamı \(64\) ise en büyüğü kaçtır? \(\star\)
\item Bir sayının \(3\) fazlasının yarısı, sayının \(3\) eksiğine eşittir. Sayı kaçtır? \(\star\)
\item Baba \(40\), oğlu \(12\) yaşındadır. Kaç yıl sonra babanın yaşı oğlunun yaşının \(2\) katı olur? \(\star\)
\item Bir kalem ile bir defter toplam \(50\) TL'dir; defter kalemden \(14\) TL pahalıdır. Kalem kaç TL'dir?
\end{enumerate}
\column{0.32\textwidth}
\pause
\textbf{Cevaplar:}
\begin{enumerate}
\item \(16\)\\[8pt]
\item \(19\) \ (\(13,15,17,19\))\\[8pt]
\item \(9\)\\[8pt]
\item \(16\) yıl sonra\\[8pt]
\item \(18\) TL
\end{enumerate}
\end{columns}
\vspace{4pt}{\footnotesize \(\star\) işaretli soruların adım adım çözümü sonraki slaytlardadır.}
\end{frame}

% ---- Ders 3 alıştırma çözümleri ----
\soru{Ders 3 -- Alıştırma Çözümü (Soru 2)}
{Ardışık dört tek sayının toplamı \(64\) ise en büyüğü kaçtır?}
{Ardışık tek sayılar ikişer artar: \(x,\ x+2,\ x+4,\ x+6\).
\[4x+12=64\ \Rightarrow\ 4x=52\ \Rightarrow\ x=13\]
Sayılar: \(13,15,17,19\) \ (toplam \(64\) \checkmark) \(\Rightarrow\) en büyüğü \(19\).}

\soru{Ders 3 -- Alıştırma Çözümü (Soru 3)}
{Bir sayının \(3\) fazlasının yarısı, sayının \(3\) eksiğine eşittir. Sayı kaçtır?}
{Sayı \(x\): \ ``\(3\) fazlasının yarısı'' \(=\dfrac{x+3}{2}\), \ ``\(3\) eksiği'' \(=x-3\).
\[\frac{x+3}{2}=x-3\ \Rightarrow\ x+3=2x-6\ \Rightarrow\ x=9\]
Sağlama: \(\dfrac{12}{2}=6=9-3\) \checkmark}

\soru{Ders 3 -- Alıştırma Çözümü (Soru 4)}
{Baba \(40\), oğlu \(12\) yaşındadır. Kaç yıl sonra babanın yaşı oğlunun yaşının \(2\) katı olur?}
{\(t\) yıl sonra: baba \(40+t\), oğul \(12+t\) (ikisi de \(t\) yaş büyür).
\[40+t=2(12+t)\ \Rightarrow\ 40+t=24+2t\ \Rightarrow\ t=16\]
Sağlama: \(16\) yıl sonra baba \(56\), oğul \(28\); \ \(56=2\cdot 28\) \checkmark}

% =====================================================
\begin{frame}{Özet}
\begin{columns}[t]
\column{0.32\textwidth}
\begin{block}{Ders 1}
\small
\begin{itemize}
\item \(ax+b=0,\ a\neq0\): \(x=-\dfrac{b}{a}\)
\item İki tarafa aynı işlem
\item Karşıya geçen işaret değiştirir
\item \(0\)'a bölme yok
\item Her zaman sağla
\end{itemize}
\end{block}
\column{0.32\textwidth}
\begin{block}{Ders 2}
\small
\begin{itemize}
\item Kesirli: EKOK ile çarp
\item Ondalıklı: \(10^{k}\) ile çarp
\item \(a=0,b=0\): \(\text{ÇK}=\mathbb{R}\)
\item \(a=0,b\neq0\): \(\text{ÇK}=\emptyset\)
\item Payda \(\neq 0\) koşulu
\end{itemize}
\end{block}
\column{0.32\textwidth}
\begin{block}{Ders 3}
\small
\begin{itemize}
\item \(3x+5\neq 3(x+5)\)
\item Ardışık: \(x, x+1, \dots\)
\item Çift/tek: \(x, x+2, \dots\)
\item Yaş: herkese aynı süre
\item Yol \(=\) hız \(\times\) zaman
\end{itemize}
\end{block}
\end{columns}
\vspace{6pt}
\begin{alertblock}{Gelecek hafta}
Birinci dereceden iki bilinmeyenli denklemler ve denklem sistemleri
\end{alertblock}
\end{frame}

\begin{frame}
\centering\Large İyi çalışmalar
\end{frame}

\end{document}
